% 7
\begin{frame}[t]{评估指标与解释边界}
\vspace{.15cm}
\fontsize{8.7}{12}\selectfont\renewcommand{\arraystretch}{1.16}
\begin{tabularx}{\textwidth}{@{}P{2.35cm}Y Y@{}}
\toprule\textbf{要回答的问题}&\textbf{拟采用的测量}&\textbf{解释边界}\\\midrule
变化量是否准确 & 预测与真实的 $\Delta\bar{x}=\bar{x}_{\mathrm{pert}}-\bar{x}_{\mathrm{ctrl}}$ 之间的 RMSE / MAE。 & 全基因相关性易被共享响应抬高；按高变基因与差异基因分层报告。\\
是否分得清扰动 & Systema 类扰动区分：真实扰动在候选扰动中的排名或 centroid accuracy。 & 依赖候选扰动数量与相似度；需与随机、均值基线比较。\\
差异基因是否找对 & 预测与真实差异基因集合的重合率与方向一致性。 & 受差异检验方法和阈值影响；固定检验流程后再比较。\\
是否超过基线 & 同划分、同指标下与对照均值、平均效应、线性 / 加性基线的差距。 & 基线选择决定结论强弱；报告全部基线而非最弱者。\\
跨背景是否泛化 & 按细胞背景留出；目标背景只提供对照数据。 & 预训练模型可能见过公开数据；核对数据暴露并如实说明。\\\bottomrule
\end{tabularx}
\source{拟议指标组合；来源 [10,11,13]。$\Delta\bar{x}$ 与区分指标的操作定义见附录，非所有论文共用的官方指标。}
\end{frame}
